% ============================================================================
% TWIN PRIMES AND SIEGEL ZEROS: A FULCRUM
% v0.1 drafted 2026-07-19 (Dichotomy Day) per the noon council rulings
% (docs/exploration/council-brief-noon.md, D1-D6 + trust-surface + discovery).
% Target: publication approval -> arXiv (math.NT primary, cs.LO cross) -> Pi.
% Convention: every stated result carries its Lean name via \leanname{};
% Appendix A maps statement -> declaration -> axiom audit.
% Voice: problem-first; every strong adjective paid by a number; every
% conditional named at the statement. No local LaTeX toolchain on the build
% machine -- compile checks happen author-side/arXiv.
% ============================================================================
\documentclass[11pt]{amsart}
\usepackage[margin=1.1in]{geometry}
\usepackage{amsmath,amssymb,amsthm}
\usepackage{booktabs}
\usepackage[colorlinks=true,linkcolor=blue!60!black,citecolor=blue!60!black,urlcolor=blue!60!black]{hyperref}
\usepackage{microtype}
\usepackage{xcolor}

\newcommand{\leanname}[1]{\par\noindent{\small\texttt{#1}}}
\newcommand{\leaninline}[1]{{\small\texttt{#1}}}

% --- Appendix A transcribes Lean source. `alltt' is base LaTeX's verbatim-like
% environment that keeps \ { } active, which is what lets the Unicode below be
% mapped to math glyphs; the mappings need a LaTeX kernel of 2018 or later
% (UTF-8 input by default), or on an older one add \usepackage[utf8]{inputenc}.
% \lchar sets each glyph in one typewriter-character box (measured from the
% current font, so the listings stay column-aligned under any tt family).
% Every non-ASCII character occurring in Appendix A is declared here; the list
% is exhaustive and was generated from the appendix source.
\usepackage{alltt}
\newcommand{\lchar}[1]{\hbox to \fontcharwd\font`\M{\hss$#1$\hss}}
\DeclareUnicodeCharacter{00AC}{\lchar{\neg}}
\DeclareUnicodeCharacter{039B}{\lchar{\Lambda}}
\DeclareUnicodeCharacter{03B1}{\lchar{\alpha}}
\DeclareUnicodeCharacter{03B2}{\lchar{\beta}}
\DeclareUnicodeCharacter{03B4}{\lchar{\delta}}
\DeclareUnicodeCharacter{03B5}{\lchar{\varepsilon}}
\DeclareUnicodeCharacter{03B8}{\lchar{\theta}}
\DeclareUnicodeCharacter{03BC}{\lchar{\mu}}
\DeclareUnicodeCharacter{03BE}{\lchar{\xi}}
\DeclareUnicodeCharacter{03C1}{\lchar{\rho}}
\DeclareUnicodeCharacter{03C7}{\lchar{\chi}}
\DeclareUnicodeCharacter{03C9}{\lchar{\omega}}
\DeclareUnicodeCharacter{2016}{\lchar{\Vert}}
\DeclareUnicodeCharacter{2080}{\lchar{{}_0}}
\DeclareUnicodeCharacter{208A}{\lchar{{}_+}}
\DeclareUnicodeCharacter{2102}{\lchar{\mathbb{C}}}
\DeclareUnicodeCharacter{2115}{\lchar{\mathbb{N}}}
\DeclareUnicodeCharacter{211A}{\lchar{\mathbb{Q}}}
\DeclareUnicodeCharacter{211D}{\lchar{\mathbb{R}}}
\DeclareUnicodeCharacter{2124}{\lchar{\mathbb{Z}}}
\DeclareUnicodeCharacter{2192}{\lchar{\to}}
\DeclareUnicodeCharacter{2200}{\lchar{\forall}}
\DeclareUnicodeCharacter{2202}{\lchar{\partial}}
\DeclareUnicodeCharacter{2203}{\lchar{\exists}}
\DeclareUnicodeCharacter{2208}{\lchar{\in}}
\DeclareUnicodeCharacter{220F}{\lchar{\prod}}
\DeclareUnicodeCharacter{2211}{\lchar{\sum}}
\DeclareUnicodeCharacter{2223}{\lchar{\mid}}
\DeclareUnicodeCharacter{2227}{\lchar{\wedge}}
\DeclareUnicodeCharacter{2228}{\lchar{\vee}}
\DeclareUnicodeCharacter{222B}{\lchar{\int}}
\DeclareUnicodeCharacter{2260}{\lchar{\neq}}
\DeclareUnicodeCharacter{2264}{\lchar{\le}}
\DeclareUnicodeCharacter{2265}{\lchar{\ge}}
\DeclareUnicodeCharacter{230A}{\lchar{\lfloor}}
\DeclareUnicodeCharacter{230B}{\lchar{\rfloor}}

\newtheorem{theorem}{Theorem}[section]
\newtheorem{corollary}[theorem]{Corollary}
\newtheorem{definition}[theorem]{Definition}
\newtheorem{proposition}[theorem]{Proposition}
\theoremstyle{remark}
\newtheorem{remark}[theorem]{Remark}

\title{Twin Primes and Siegel Zeros: a Fulcrum}
\author{Jason Hickey}
\thanks{Developed in collaboration with Claude (Anthropic's Fable~5 and
Opus models); see \S\ref{sec:method} and the repository's commit ledger,
in which every commit is co-authored. Written in a personal capacity.}
\date{July 19, 2026}

\begin{document}

\begin{abstract}
Every theorem in this paper is verified by the Lean~4 kernel over
\textsf{mathlib}, with axiom base exactly
$\{\texttt{propext}, \texttt{Classical.choice}, \texttt{Quot.sound}\}$;
nothing rests on informal argument. The mathematics---the designs, the
proofs, and their formalization---is joint work of the author and Claude
(Anthropic's Fable~5 and Opus models), carried out under the author's
direction and ratification; the commit ledger records the collaboration
line by line.

We introduce the \emph{fulcrum}: a single hypothesis---one fixed quality
constant $C^\star$, one zero per unbounded window---under which, together
with an explicitly stated engine hypothesis, Siegel zeros produce
infinitely many twin primes; the fulcrum is proved minimal for the
engine that consumes it, with the
zero's reality \emph{derived} rather than assumed. The fulcrum was not
transcribed from the literature: it was extracted from a formal corpus by
demand-audited hypothesis minimization, a mode of discovery we argue is
impossible over prose proofs. Around it we prove, kernel-checked: a
Heath-Brown--type dichotomy (twin primes, or no Siegel zeros) with its
engine-conditional horn stated explicitly; the first formalization of any
zero-free region for $\zeta$ of Littlewood strength (1922) and the first
power-saving region ($\theta = 3/4$) in any proof assistant; a
Deuring--Heilbronn repulsion contract; the exact-identity keystone of the
Heath-Brown engine, proven unconditional; the \emph{exchange-rate
wall}, a neutrality theorem over arbitrary real sign functions that prices
every transfer-based route to twins; and the \emph{parity gap} as a
mathematical object: a predicate $Z$ with a kernel-checked theorem that no
true, twin-sufficient, parity-invariant completion exists, together with a
census placing the entire corpus inside the parity-invariant cone. As a
byproduct, the logarithmic two-point Chowla statement---a theorem of
Tao---is machine-checked conditionally on a single named uniformity
hypothesis. The corpus is $324{,}724$ lines of Lean across
$751$ files; the method that produced it---the \emph{Salt method}---kept a
public ledger of every design error caught (256 numbered entries) against
zero wrong proofs, and is described as a contribution in its own right.
\end{abstract}

\maketitle

% ============================================================================
\section{Introduction}\label{sec:intro}
% ============================================================================

The twin prime conjecture is not merely unproven; the routes toward it are
walled by a structural obstruction---the parity barrier---that has never
itself been a theorem, and guarded by a ghost---the Siegel zero---that has
never been an object one could bargain with precisely. Progress of the
conditional kind (Heath-Brown's 1983 theorem that infinitely many Siegel
zeros force infinitely many twin primes \cite{HB1983}) trades in
hypotheses whose exact strength is scattered through pages of estimates:
what, \emph{precisely}, must one believe about exceptional zeros to
believe in twin primes? Prose proofs cannot answer such questions
exhaustively, because a prose proof hides its demand.

A formal proof \emph{is} its demand. This paper's central object, the
\textbf{fulcrum}, was found by exploiting exactly that: we formalized a
Heath-Brown--type engine, enumerated---mechanically, exhaustively---every
site at which the Siegel-zero hypothesis is consumed and at what strength,
and then weakened the hypothesis axis by axis to the audited demand, with
the Lean kernel refereeing each weakening instantly. The fixed point of
that process is the weakest twin-producing Siegel hypothesis yet
formulated (Definition~\ref{def:fulcrum}); the process itself,
\emph{demand-audited hypothesis minimization}, is to our knowledge a new
mode of mathematical discovery, and it is the sense in which this paper's
formalization is an instrument rather than a notary.

The results divide into four movements. First, the \emph{banners}
(\S\ref{sec:banners}): unconditional analytic theorems formalized for the
first time, including a Littlewood-strength zero-free region for $\zeta$
and the first power-saving region in any proof assistant. Second, the
\emph{fulcrum and the dichotomy} (\S\S\ref{sec:fulcrum}--\ref{sec:keystone}):
the minimized hypothesis, the kernel-checked dichotomy it induces, and the
exact-identity keystone of the engine, now proven unconditional. Third,
the \emph{negative space}
(\S\S\ref{sec:walls}--\ref{sec:parity}): the exchange-rate wall and the
parity gap, which together map---as theorems, not folklore---why the
obvious routes die. Fourth, the \emph{method} (\S\ref{sec:method}): the
Salt method, an adversarial division of mathematical labor among human
direction, frontier-model design, fleet-model execution and refutation,
and kernel verification, whose public error \emph{ledger} --- an
append-only record in the repository of every design error caught,
numbered and repaired, whose entries we cite as (ledger \#$N$) --- stands
at 256 numbered catches against zero wrong proofs and is offered as evidence that the standard of the first
sentence of the abstract is achievable at research scale and research
speed. All of it was produced by the Salt method (\S\ref{sec:method}).

Throughout, conditionality is a first-class citizen: every conditional
statement names its exact residual hypothesis, and
Table~\ref{tab:conditional} lists all of them in one place. The paper
makes no unconditional claim about twin primes.

% ============================================================================
\section{The corpus and the referee}\label{sec:trust}
% ============================================================================

The corpus is $324{,}724$ lines of Lean~4 over \textsf{mathlib}
\cite{mathlib}, in $751$ files, built by $1034$ commits; a full build
kernel-checks roughly $9{,}400$ compilation jobs. The skeptic's question
---\emph{how do we really know?}---splits into what is mechanically
excluded and what is irreducibly a human question, and we answer it in
that order.

\subsection{Mechanically excluded}
A \texttt{sorry} anywhere in a dependency tree---ours or
\textsf{mathlib}'s---introduces the axiom \texttt{sorryAx}, which
propagates transitively and surfaces in the axiom closure of every
downstream theorem. Our \texttt{\#audit\_axioms} blocks print that closure
for every headline statement: it is exactly
$\{\texttt{propext}, \texttt{Classical.choice}, \texttt{Quot.sound}\}$,
the three classical axioms of \textsf{mathlib} practice. A hidden gap
cannot hide. Tactic bugs cannot certify falsehoods: tactics emit proof
terms that the kernel re-checks independently, so the trust surface is the
kernel alone, a deliberately small core; we additionally forbid
\texttt{native\_decide} (which would trust the compiler, visibly, as a
fourth axiom). The kernel itself is the one residual mechanical trust,
and it is checkable: before submission every declaration is replayed
through \texttt{leanchecker}, the toolchain's standalone proof checker,
independently of the elaborator and build cache that produced it. The
replay does not remove the kernel from the trusted base---checker and
kernel share an implementation---but it excludes elaborator and cache
faults, which is where the practical risk lives. One consequence for the reader: every
numbered statement in this paper---theorem, proposition,
corollary---is kernel-checked; the labels grade importance, never
epistemic status.

\subsection{The specification gap}
What no formal system can check is whether our \emph{statements} mean what
our prose claims. This gap is irreducible---it is the map--territory
boundary---so we engineer against it. Definitions bottom out in
\textsf{mathlib}'s communal objects (\texttt{Nat.Prime},
\texttt{DirichletCharacter}, \texttt{LFunction}), vetted by thousands of
independent uses. Vacuity is an attacked property, not an assumed one:
adversarial review passes carry mandatory degenerate-instance checks, and
where a definition provably trivializes outside its intended window, the
degeneracy is \emph{stated in Lean} beside it (see the grade guard of
\S\ref{sec:parity}). Non-vacuity is proved by witness: the fulcrum's
zero-reality is derived from our own zero-free region
(Theorem~\ref{thm:fulcrumreal}); the parity census instantiates its cone
with ten corpus theorems; classical corollaries are re-derived from our
statements, which a misstated premise would break. Finally,
Appendix~\ref{app:index} renders every headline statement in Lean and in
prose side by side, reducing the entire correctness question of the corpus
to a finite, public, refereeable translation audit. ``Absolutely'' is not
on offer anywhere in mathematics; what is on offer here is a doubt surface
consisting of one small kernel, three axioms, and a one-page table.

% ============================================================================
\section{The four banners}\label{sec:banners}
% ============================================================================

Four unconditional theorems anchor the corpus; none had been formalized
before, and two had never been formalized at any strength in any proof
assistant.

\begin{theorem}[Littlewood-strength zero-free region, 1922]\label{thm:litt}
There exist effective constants such that $\zeta(s) \neq 0$ in a region of
Littlewood type
$\sigma \ge 1 - c\,\frac{\log\log |t|}{\log |t|}$ for $|t|$ large.
\leanname{Salt.Vk.zeta\_zero\_free\_region\_littlewood}
\end{theorem}

\begin{theorem}[The first power-saving region]\label{thm:pow}
There exist effective constants such that $\zeta(s) \neq 0$ for
$\sigma \ge 1 - c\,(\log|t|)^{-3/4}(\log\log|t|)^{-O(1)}$, $|t|$ large:
a Vinogradov--Korobov--type region at exponent $\theta = 3/4$.
\leanname{Salt.Vk.zeta\_zero\_free\_region\_pow}
\end{theorem}

To our knowledge Theorem~\ref{thm:pow} is the first zero-free region with
a power saving in the exponent verified in any proof assistant, and
Theorem~\ref{thm:litt} the first formalization of the 1922 Littlewood
region; the strongest previously formalized regions were of classical
de la Vall\'ee Poussin shape. The proofs formalize the Vinogradov mean
value theorem (\leaninline{Salt.Vmvt.vmvt}) and route it through an
exponential-sum-to-growth bridge (\leaninline{Salt.Vk.vk\_block\_core}),
with an explicit strip family at absolute constant $C = 4096$.

\begin{theorem}[Deuring--Heilbronn repulsion, contract form]\label{thm:rep}
There are explicit constants $b = 680$, $k = 14$ and $c > 0$ such that: for
every primitive quadratic character $\chi \bmod q$, $q \ge 2$, every real
zero $\beta_0 \in (\tfrac12, 1)$ of $L(s,\chi)$, and every non-real zero
$\rho$ of $L(s,\chi)$ with
$\tfrac{16}{17} \le \mathrm{Re}\,\rho \le \beta_0$ --- at \emph{every}
height ---
\[
1 - \beta_0 \;\ge\;
  \frac{c\,\bigl(q(|\mathrm{Im}\,\rho| + 2)\bigr)^{-b(1 - \mathrm{Re}\,\rho)}}
       {\bigl(\log\bigl(q(|\mathrm{Im}\,\rho| + 2)\bigr) + 2\bigr)^{k}}.
\]
The closer the real zero sits to $1$, the farther every companion zero is
pushed from the $1$-line: the Deuring--Heilbronn repulsion phenomenon, in
ordered two-zero contract form with every constant explicit.
\leanname{Salt.SW.dh\_repulsion\_tall}
\end{theorem}

\begin{remark}[The constant, and what sharpening it buys]\label{rem:repc}
The height restriction $|\mathrm{Im}\,\rho| \le 1$ is not needed: it was
removed from the contract outright rather than widened to a parameter, and
the height-bounded twin survives beside it at the same $b$ and $k$
(\leaninline{Salt.SW.dh\_repulsion\_ordered}). The constant $c$ is not a
single symbol but an explicit minimum of ten arms, each the price of one row
of the underlying estimate; in the unrestricted form every arm is free of the
auxiliary $\zeta$-growth constant, so the tower prices exactly:
\[
\log(1/c) \;=\; 86.2267 \qquad \text{at } c_0 = 1/126848,
\]
$c_0$ being the corpus's own zero-free-region constant, and the binding arm
being $(c_0/32)^{17/3}$ --- that is, the repulsion constant is now limited by
our zero-free region rather than by the estimate chain that consumes it. The
figure was $\log(1/c) = 631.5764$ when the contract first landed; two
sharpening waves (replacing a uniform exponent by each row's own H\"older
floor, and retiring a $2^{-250}$ arm) brought it to the value above
\emph{with no parameter changed and no statement outside the two contract
files touched}. We record what this does not buy: the dichotomy of
\S\ref{sec:dichotomy} is stated for every $C > 0$, so no threshold in this
paper moves with the size of $c$; the gain is explicit-constant hygiene, and
it matters only where a downstream hypothesis packet must be shown
non-empty.
\end{remark}

\begin{theorem}[Vinogradov mean value theorem]\label{thm:vmvt}
The mean value estimate for Weyl sums in the range needed by
Theorem~\ref{thm:pow}, with explicit constants.
\leanname{Salt.Vmvt.vmvt}
\end{theorem}

Around these sit an unconditional Siegel--Walfisz theorem, a
Bombieri--Vinogradov theorem, a machine-checked Chen's theorem with the
unconditional corollary $\Omega(p(p+2)) \le 3$ for infinitely many primes
$p$ (\leaninline{Salt.Fulcrum.chen\_omega\_prod\_le\_three}), and the
unconditional $\Omega(n(n+2)) \le 20$ infinitely often
(\leaninline{Salt.BrunLower.twin\_almost\_prime}); these are inherited
infrastructure from the same program and are indexed in
Appendix~\ref{app:index}.

% ============================================================================
\section{The fulcrum}\label{sec:fulcrum}
% ============================================================================

\subsection{The discovery}
The Heath-Brown engine \cite{HB1983} consumes a Siegel-zero hypothesis at many
sites: as a lower bound forcing $L$-values small, as a localization of one
zero, as a quality bound entering a sieve level, as a reality assumption.
In prose, the totality of this demand is not enumerable; in the formal
corpus it is a finite list of typed consumption sites, each with an exact
strength. The Fulcrum Hunt (three adversarial passes over the corpus,
recorded in the repository's exploration ledger) audited every site,
then weakened the hypothesis along each axis until the audited demand was
met exactly---each candidate weakening compiled or refuted by the kernel
within minutes. Three collapses survived: the quality constant collapsed
from a universally quantified family to a \emph{single fixed} $C^\star$;
the zero's \emph{reality} moved from hypothesis to conclusion (derived
from our own zero-free region); and the \emph{quantity} collapsed to one
zero per unbounded window---no asymptotic family, no density, no
infinitude of characters beyond what one zero per window forces.

\begin{definition}[The fulcrum]\label{def:fulcrum}
For $C > 0$, \emph{FulcrumQualityMin} $C$ asserts: for every natural
number $Q$ there exist a modulus $q > Q$, a primitive quadratic
character $\chi \bmod q$, and a zero $\rho$ of $L(\cdot,\chi)$ with
$\|1-\rho\| \cdot C \log q \le 1$.
The \emph{fulcrum hypothesis} is
\[
F \;:=\; \text{FulcrumQualityMin } C^\star
\]
for the single explicit constant $C^\star = \max(C^{(1)}, 2/c_0)$, where
$c_0 = 1/126848$ is the zero-free-region constant of the corpus,
kernel-certified at that numeral
(\leaninline{zero\_free\_region\_all\_numeral}); the $2/c_0$ arm of
$C^\star$ is exactly $253696$, and the fulcrum's reality theorem is
instantiated at that bare threshold
(\leaninline{fulcrum\_zero\_real\_numeral}).
\leanname{Salt.Fulcrum.FulcrumQualityMin}
\end{definition}

\begin{theorem}[Reality is derived]\label{thm:fulcrumreal}
Any zero witnessing the fulcrum at quality $C^\star$ is real, and lies in
the exceptional position of Siegel-zero theory.
\leanname{Salt.Fulcrum.fulcrum\_zero\_real}
\end{theorem}

\begin{remark}
The threshold $Q$ is a bare cursor: the quantifier
$\forall Q\,\exists q > Q$ asserts exactly that the qualifying moduli
form an unbounded set --- no density, frequency, or growth condition is
imposed, and each instance demands a single witness. In particular, one
Siegel zero is one bounded window: finding a single exceptional zero
would refute GRH but would \emph{not} satisfy the fulcrum, which demands
a witness beyond every $Q$. The fulcrum carries no asymptotic content
per window; this is the weakest shape for which the engine of
\S\ref{sec:keystone} can run, and the demand audit certifies that
\emph{this} engine leaves nothing weaker on the table.
\end{remark}

\begin{remark}[The two arms of $C^\star$]\label{rem:arms}
Only the second arm of $C^\star = \max(C^{(1)}, 2/c_0)$ is a numeral. The
first, $C^{(1)} = \exp\exp\{2A/(\mathfrak{S}\,C(\alpha))\}$ in the
notation of \cite{HB1983} (its Corollary~2, with $A$ the effective
constant of its Corollary~1), is the threshold at which that engine's
error budget closes. It is a double exponential: an inner exponent as
small as $2A/(\mathfrak{S}\,C(\alpha)) = 10$ already places $C^{(1)}$
near $e^{22026}$, against the reality arm's $2/c_0 = 253696$. So
$C^{(1)}$ decides the maximum, and the constant this paper exhibits as a
numeral is the reality floor---the quality at or above which
Theorem~\ref{thm:fulcrumreal} derives that the witnessing zero is
real---not the quality at which any engine is expected to run.

No statement here depends on the size of that gap.
Theorem~\ref{thm:dich} quantifies $C$ away: it is stated for every
$C > 0$ and consumes the engine implication at that same $C$, so an
engine supplying $h_{\mathrm{Engine}}$ at $C^{(1)}$ instantiates the
dichotomy at $C^{(1)}$, with the same conclusion. The reality derivation
runs in the same direction---its hypothesis is $C \ge 2/c_0$, so it
applies a fortiori at any threshold above the numeral. We report
$253696$ because it is the arm the corpus certifies
(\leaninline{fulcrum\_zero\_real\_numeral}), not as a prediction of where
$h_{\mathrm{Engine}}$ will be supplied.
\end{remark}

% ============================================================================
\section{The dichotomy}\label{sec:dichotomy}
% ============================================================================

Here \emph{TwinPrimeConjecture} is the expected statement, in the same
unbounded-threshold form as the fulcrum's own quantifier:
\[
\forall\, n\ \exists\, p \ge n:\ p \text{ and } p+2 \text{ prime}
\]
(\leaninline{Salt.Basic: TwinPrimeConjecture}) --- the dichotomy trades
one unboundedness for another.

\begin{theorem}[The dichotomy, with the engine hypothesis
explicit]\label{thm:dich}
Unconditionally: if the fulcrum fails ($\lnot F$), there are no Siegel
zeros in the precise sense of the corpus's \leaninline{NoSiegelZeros}
(\leaninline{Salt.Fulcrum.not\_fulcrum\_implies\_noSiegelZeros}).
Consequently, for any $C > 0$, given the engine implication
$h_{\mathrm{Engine}} : \text{FulcrumQualityMin } C \to
\text{TwinPrimeConjecture}$, we obtain
\[
\text{TwinPrimeConjecture} \ \lor\ \text{NoSiegelZeros}.
\]
\leanname{Salt.Fulcrum.fulcrum\_dichotomy}
\end{theorem}

Here \emph{NoSiegelZeros} is precisely the statement
\begin{gather*}
\exists\, c > 0:\ \forall\, q > 1,\ \forall\, \chi \bmod q
\text{ primitive quadratic},\ \forall\, \beta \in \mathbb{R}:\\
L(\beta, \chi) = 0 \ \wedge\ \beta < 1
\ \Longrightarrow\ \beta \le 1 - \frac{c}{\log q},
\end{gather*}
i.e.\ some single effective constant pushes every real zero of every
primitive quadratic character---all moduli, no exceptions---away from
$s = 1$ at the scale $1/\log q$. Since ``Siegel zero'' is inherently a
threshold-relative notion, this $\exists c$ form is the exact formal
content of the phrase ``there are no Siegel zeros'' as used in the
literature (it is Tao's formulation); it is far weaker than GRH, and it
speaks only of real zeros of quadratic characters---the classical
exceptional configuration, which is the only one that can occur.

The second horn is thus the Heath-Brown dichotomy, machine-checked and
strengthened: strengthened because the hypothesis feeding its twin horn is
the minimized fulcrum rather than the 1983 hypothesis, and because the
$\lnot F$ horn is unconditional and effective. The engine implication
$h_{\mathrm{Engine}}$ is the one conditional input; \S\ref{sec:keystone}
proves its exact-identity keystone unconditionally, leaving the remaining
Heath-Brown assembly as the named campaign \textsc{hb-engine}. We do not
claim the unconditional dichotomy; we exhibit exactly what
remains.

% ============================================================================
\section{The keystone}\label{sec:keystone}
% ============================================================================

The objects: for a primitive quadratic $\chi \bmod q$, the twisted von
Mangoldt function is
\[
\tilde\Lambda(n) \;:=\; \sum_{d \mid n} \mu^2(d)\,\chi(d)\,\log(n/d),
\]
a majorant-with-sign of $\Lambda$ that agrees with it when $\chi$ is
$\lambda$-like on the divisors of $n$; over a finite window $A$,
\[
S_1(A) := \sum_{n \in A} \Lambda(n)\Lambda(n+2), \qquad
S_2(\chi, A) := \sum_{n \in A} \tilde\Lambda(n)\tilde\Lambda(n+2),
\]
and the \emph{pretense sum} --- the price currency of the whole chapter ---
is
\[
\mathrm{PS}_\chi(N) \;:=\; \sum_{\substack{p \le N \\ \chi(p) = 1}}
\frac{\log p}{p},
\]
which is $O(1)$ exactly when $\chi$ pretends to be the Liouville pattern
on primes. The window $W = W(\chi, z, x)$ is the set of $n \in (x, 2x]$
with $n(n+2)$ coprime to $q$ and to every prime $p < z$ with
$\chi(p) \ne -1$: the exact sifted support on which
$\tilde\Lambda \ge \Lambda \ge 0$ termwise, so $S_1 \le S_2$ and the
engine's content is entirely in bounding the overshoot $S_2 - S_1$.

The engine's hardest block is exactly that second-moment comparison
$S_2 - S_1$---the site of Heath-Brown's \cite{HB1983} Lemma~2. Formalization exposed that the classical route's
$\tau$-crude majorant is \emph{provably} insufficient at the needed
generality (its concentration pattern is unexcludable by pattern-blind
tools; ledger entry \#251), so the corpus proves the required conclusion
directly on an exact identity: the overshoot
$(\tilde\Lambda-\Lambda)(n)\tilde\Lambda(n+2) +
\Lambda(n)(\tilde\Lambda-\Lambda)(n+2)$ is summed exactly, its support is
classified, and eleven family estimates (each with explicit absolute
constants; nine landed at first attempt) price every class.

\begin{theorem}[The master estimate, unconditional]\label{thm:master}
With $W$ the sifted window at parameters $(z,x)$ in the engine's regime:
\[
S_2(\chi, W) - S_1(W) \;\le\; 2^{31}\Bigl(\frac{x}{z_0}
+ \frac{x}{\log x}\,e^{5 z_0}\,\mathrm{PS}_\chi
+ e^{2 z_0}\bigl(x z^{-1/8} + x^{9/10}\bigr) L^3\Bigr),
\]
where $z_0 = \log(2x+2)/\log z$, $L = \log(2x+2)$, and
$\mathrm{PS}_\chi$ is the pretense sum of $\chi$.
\leanname{Salt.HB.hb\_l2c\_master\_unconditional}
\end{theorem}

What briefly appeared to be a residual---a $\chi$-weighted sieve count
(the roles-swapped mirror family, where the character-positive prime is a
cofactor rather than a modulus, together with one boundary family of the
same shape at the engine's witness scale)---was \emph{proven
unnecessary}: both counts route through the already-established
pair-count engine \emph{character-blind}, landing in the leading budget
row, so the estimate holds from its four explicit hypotheses alone. The
cover of the eleven families is itself a kernel-checked theorem---no case
can have been forgotten---and the three cover gaps found during execution
were all found \emph{by the machine layers below the kernel} and repaired
at the statement level before any proof was attempted on a wrong
statement (ledger \#245, \#246, \#252). The regime bookkeeping produced
one of the campaign's sharpest catches, an amendment refuted by its own
executor at the asymptotic corner (ledger \#253); at the downstream
witnesses the full hypothesis packet is discharged by explicit constants
(\leaninline{Salt.HB.glue\_master}).

% ============================================================================
\section{The exchange-rate wall}\label{sec:walls}
% ============================================================================

The Fulcrum Hunt's second product is negative space made precise: four
recorded refutations (the \emph{wall series}, banked in the repository's
exploration ledger) of the natural routes by which Siegel-zero
information might be converted into twin-prime information. Three are
ledger-grade records; the fourth is a theorem series, formalized here.

The transfer at the heart of every such route exchanges the twin sum
$S_1$ for a twisted sum $S_2(f)$ against a multiplicative sign function
$f$. The wall states: this exchange is \emph{never free, and its price is
exactly a pretense sum}---and it says so for arbitrary real sign
functions, a class strictly wider than Dirichlet characters.

\begin{definition}[Sign functions]\label{def:sign}
$f : \mathbb{N} \to \mathbb{R}$ is a \emph{sign function} if it is
totally multiplicative with $f(1)=1$, $|f| \le 1$, and $f(p) = \pm 1$ on
primes.
\leanname{Salt.HB.IsSignFunction}
\end{definition}

\begin{theorem}[Neutrality and the rate]\label{thm:wall}
For every sign function $f$: $S_1 \le S_2(f)$ on the generalized window
(unconditionally), and the difference obeys the budget
\[
S_2(f) - S_1 \;\le\; \mathrm{EngineBound}(C_{\mathrm{main}}, z, x,
\mathrm{PS}_f)
\]
whenever the overshoot budget holds---the three-row bound of
Theorem~\ref{thm:master} with the character's pretense sum replaced by
$f$'s.
\leanname{Salt.HB.S1\_le\_S2Gen, Salt.HB.neutrality\_rate}
\end{theorem}

\begin{theorem}[The neutrality point]\label{thm:liouville}
The Liouville function is the exact fixed point of the exchange:
$\tilde\Lambda_\lambda = \Lambda$ identically, $S_2(\lambda) = S_1$, the
overshoot vanishes pointwise, and $\mathrm{PS}_\lambda = 0$.
\leanname{Salt.HB.LamTildeGen\_lamR\_eq\_vonMangoldt,
S2Gen\_lamR\_eq\_S1, overshootExactGen\_lamR,\\
PretenseSumGen\_lamR\_eq\_zero}
\end{theorem}

Read together: any route that would extract twins from a transfer must pay
in pretense, the currency $\mathrm{PS}_f$; the one sign function with
$\mathrm{PS} = 0$ is the one at which the transfer returns exactly the
twin sum it was fed. The wall is falsifiable---exhibit a sign function
with bounded pretense sum and an $S_2$-evaluation that does not route
through $S_1$---and it is stated over the full sign-function cone
precisely so that its scope is a theorem rather than an impression. The
re-typing of the entire transfer chain from Dirichlet characters to sign
functions (59 declarations) required deleting every coprimality thread;
that the deletion is sound is itself part of the formal content: the
character was never load-bearing.

% ============================================================================
\section{The parity gap, as an object}\label{sec:parity}
% ============================================================================

The parity barrier has lived in the literature as folklore with
heuristics. The corpus makes it an object with a theorem.

\begin{definition}[Completions, parity invariance, and $Z$]\label{def:z}
Let $\rho(d) = \#\{0 \le r < d : d \mid r(r+2)\}$ be the true twin local
density. For a weighting $a : \mathbb{N} \to \mathbb{R}$, the type-I sum
at level $d$ and window $x$ and its error at exponent $\theta$ are
\[
A_d(a; x) := \sum_{\substack{1 \le n \le x \\ d \mid n(n+2)}} a(n),
\qquad
E_\theta(a; x) := \sum_{d \le x^{\theta}}
  \Bigl|\,A_d(a;x) - \frac{\rho(d)}{d}\,x\,\Bigr|.
\]
The \emph{completion class} at grade $(\theta, A_0)$ is
\[
\mathcal{C}(\theta, A_0) \;:=\; \Bigl\{\, a \ge 0 \;:\; \exists\, C > 0,\
\forall x \ge 2,\ E_\theta(a; x) \le \frac{C\,x}{(\log x)^{A_0}} \Bigr\}
\]
--- the nonnegative weightings indistinguishable from the truth by pure
type-I information at level $x^\theta$, quality $(\log x)^{-A_0}$. (The
local density $\rho$ is \emph{pinned}, not quantified: completions must
carry the true twin main term. Type-II bilinear information is
structurally inexpressible in the class, not merely excluded.) A
predicate $E$ on weightings is \emph{parity-invariant} at grade
$(\theta, A_0)$ if it holds for \emph{every} completion:
$\mathrm{ParityInv}(E) :\iff \forall a \in \mathcal{C}(\theta, A_0),\ E(a)$
--- semantic invariance over the model class, quantifying over the
mechanism rather than any method. With the twin mass
$T(a; x) := \sum_{n \le x,\ n,\,n+2 \text{ prime}} a(n)$, the predicate
$E$ is \emph{twin-sufficient} if for every $a \in \mathcal{C}(\theta,
A_0)$ satisfying $E$, the mass $T(a; \cdot)$ is unbounded. Writing
$\mathbf{1}$ for the constant weighting $1$ (the true sequence), the gap
statement is
\[
Z(\theta, A_0) \;:\iff\; \exists\, E \ \text{twin-sufficient with}\
E(\mathbf{1}).
\]
No Dirichlet character, $L$-function, or exceptional-zero oracle appears
in any of these inputs; in the formal corpus this is enforced
syntactically---the defining file imports nothing from the Siegel-zero or
sieve tracks, checkable by its import list alone.
\leanname{Salt.Parity.Z, Salt.Parity.ParityInv (Salt/Parity/Z.lean)}
\end{definition}

\begin{theorem}[The gap theorem]\label{thm:gap}
For $0 < \theta < \tfrac12$ and the certified $A_0$ range, no predicate
$E$ is simultaneously true ($E(\mathbf{1})$), twin-sufficient, and
parity-invariant. The witness is the twin-free weighting
$a^-(n) = 0$ if $n, n+2$ are both prime and $1$ otherwise: a
Brun-grade argument (using the corpus's explicit
$\pi_2(x) \ll x/\log^2 x$ bound) places $a^- \in
\mathcal{C}(\theta, A_0)$, parity-invariance carries $E$ from the truth
to $a^-$, twin-sufficiency then forces $T(a^-;\cdot)$ unbounded --- but
$T(a^-; x) = 0$ identically.
\leanname{Salt.Parity.sufficient\_true\_not\_parityInv}
\end{theorem}

\begin{theorem}[$Z$ is exactly the demand]\label{thm:ziff}
Over the certified grade window ($0 < \theta < 1/2$, $A_0$ in the stated
range): $Z(\theta,A_0) \iff$ the twin prime conjecture.
\leanname{Salt.Parity.Z\_implies\_TPC, Salt.Parity.TPC\_implies\_Z}
\end{theorem}

\begin{proposition}[The census]\label{prop:census}
Ten headline theorems of the corpus lie inside the parity-invariant
cone: the transfer's unconditional leg $S_1 \le S_2$; the $k=2$ Maynard
bound $J_1 + J_2 \le 2\log 2 \cdot I_2$ and its $\theta$-parameterised
gate failure; the impossibility that any continuous weight crosses the
twin gate; the equivalence of NoSiegelZeros with the failure of
infinitely many Siegel zeros; Chen's theorem (infinitely many primes $p$
with $p+2$ a $P_2$, i.e.\ having at most two prime factors counted
with multiplicity---the standard sieve notation, in which the ``at
most'' is exactly the parity barrier's residue); the twin almost-prime
bound $\Omega(n(n+2)) \le 20$
infinitely often; Brun's theorem (the twin reciprocal sum converges);
the Selberg twin counting bound $\pi_2(x) = O(x/\log^2 x)$; and
Chen's second theorem (every sufficiently large even $N$ is $p + q$
with $p$ prime and $q$ a $P_2$).
\leanname{Salt/Parity/Instances.lean (ten declarations, one per
theorem)}
\end{proposition}

The three results triangulate the frontier exactly: everything the corpus
proves lives inside the cone (Proposition~\ref{prop:census}); the twin
prime conjecture is equivalent to a demand $Z$
(Theorem~\ref{thm:ziff}); and that demand cannot be met inside the cone
(Theorem~\ref{thm:gap}). The barrier is now a boundary with coordinates.
Where $Z$ degenerates outside its certified window, the degeneracy is
itself a Lean lemma beside the definition (the grade guard,
\leaninline{Z\_trivial\_of\_not\_completion})---vacuity disclosed, not
discovered.

% ============================================================================
\section{The door-only spine}\label{sec:spine}
% ============================================================================

Independently of the Siegel-zero tracks, the corpus carries a full
formalization of the entropy-decrement route to the logarithmic two-point
Chowla statement, proved by Tao~\cite{TaoChowla}; the content here is the
machine-checked reduction, not the truth of the statement. Late in the
campaign its residual collapsed: the
prior interface (three coupled analytic witnesses) was \emph{proven
unsatisfiable} at its recorded constants---a freeze-panel finding,
confirmed four independent ways, ledger \#249---and the corrected
budget was proven through, leaving:

\begin{theorem}[Log-Chowla, door-only]\label{thm:spine}
There exist $\delta_0 > 0$ and a regime $R$ such that a single hypothesis
---the Matom\"aki--Radziwi\l\l--Tao uniformity door
$\mathrm{MRTUniformity}\,R\,\delta$ at any $0 < \delta \le \delta_0$---
implies the logarithmic two-point Chowla statement at $R$. A
$\forall$-quantified head version places the regime floor above any
prescribed threshold.
\leanname{Salt.Entropy.Chowla.log\_chowla\_two\_door\_only,
log\_chowla\_two\_budget\_head}
\end{theorem}

Here the logarithmic two-point Chowla statement at scale $x$ is the
standard one,
\[
\frac{1}{\log x} \sum_{n \le x} \frac{\lambda(n)\,\lambda(n+1)}{n}
\;\longrightarrow\; 0,
\]
rendered in the corpus as the failure of a quantified
$\varepsilon$-deviation over the regime's windows. The door is the
Matom\"aki--Radziwi\l\l--Tao uniformity property restricted to the
bounded large-spectrum frequency set $\Xi_H$ (of size $\le K$,
independent of $H$): for every window length $H$ in the regime's range
and every $\xi \in \Xi_H$,
\[
\int \Bigl\| \sum_{h \le H} \lambda(n + h)\, e\bigl(-\xi h / H\bigr)
\Bigr\| \, d\mu_{\log}(n) \;\le\; \delta H,
\]
where $\mu_{\log}$ is the logarithmic measure of the regime --- the
short-interval uniformity of $\lambda$ that
Matom\"aki--Radziwi\l\l--Tao's Proposition~2.4 provides in prose
mathematics, at finitely many frequencies only, which is all the proof
fires. The door is the exact target of the corpus's staged
Matom\"aki--Radziwi\l\l\ program (whose first stones, including the
closure of a long-flagged Euler-oscillation bridge, landed the same
day); the conjecture's entire residual is thus one named seam between
two formalized campaigns.

% ============================================================================
\section{The Salt method}\label{sec:method}
% ============================================================================

\begin{quote}
\emph{The Salt method is an adversarial division of mathematical labor:
human direction and ratification; frontier-model design under statement
freezes; fleet-model execution and independent refutation, tiered by
difficulty class; and kernel verification as the sole referee. Its
discipline---freeze before proving, refute before trusting, give up
loudly, bank everything---is designed to make the system catch errors
faster than it makes them.}\footnote{After water, salt.}
\end{quote}

Language models of differing proving ability come at sharply differing
inference cost, and most of a formalized proof is mechanical; the
division of labor exploits exactly that asymmetry. Every lemma is
classified before any proof attempt into one of four difficulty
classes---\emph{mechanical} (one tactic or a known lemma),
\emph{standard} (a routine multi-step argument), \emph{design}
(genuine proof architecture), \emph{open}---and routed to the cheapest
tier of model that can reliably discharge it: the most capable model
(here Fable~5) designs the campaigns, freezes the statements, and
classifies the components; \emph{executor agents}---independent
instances of a mid-tier model (here Claude Opus), dispatched one per
proof obligation with a self-contained brief, hereafter
\emph{executors}---prove them at their tier. Misclassification is safe rather than fatal: an
under-tiered executor is required to stop and report---give up early,
loudly---and the obligation re-routes upward carrying the map of where
it broke, so a grading error costs a bounded retry, never a wrong
proof. Statements are frozen by design panels whose members include
adversarial refuters with a standing default of \emph{refuted}; executors
are forbidden to alter a frozen statement---their duty on discovering one
unprovable is to stop and report, loudly; and the house that rules on such
reports is itself subject to the same audit (twice in the campaign, an
executor refuted the supervising layer's own amendment, and the record
says so: ledger \#248, \#250).

The evidence that this discipline works is the ledger: 256 numbered
catches over the campaign against \emph{zero} wrong proofs built on any of
them. The day this paper was drafted supplies the case study: five
statement-layer faults surfaced in one morning---three cover-taxonomy
gaps in a frozen design, one refutation of a supervising amendment, one
unsatisfiable recorded interface---and every one was caught by the layers
below the kernel, ruled on at the statement level, and repaired before a
single wrong proof was attempted. A companion paper (\emph{The Salt
Method}) will give the full treatment; the repository's ledgers are the
raw artifact both papers cite.

We record one measurement with a claim attached: no design error was
first discovered by the kernel. Proof attempts failed and were repaired
as routine, but every one of the ledger's catches was made by the
adversarial layers above the kernel, at the statement level, before a
wrong proof was built on it---and no landed statement was later
retracted. The kernel confirmed; the layers discovered. A verification standard is achievable at research
speed precisely because correctness is enforced \emph{above} the kernel
and only certified at it.

% ============================================================================
\section{The program forward}\label{sec:forward}
% ============================================================================

Three named campaigns continue from the corpus's current edge, each with
its target stated formally in advance: \textsc{hb-engine} (the remaining
distance from Theorem~\ref{thm:master}, now unconditional, to
$h_{\mathrm{Engine}}$); the Matom\"aki--Radziwi\l\l\ program toward the
door of Theorem~\ref{thm:spine}; and the explicit-constants floor
(Siegel--Walfisz through Chen with every constant explicit). The
conditional table below is the paper's exact boundary.

The first of the three has since acquired coordinates, which we state
because ``what remains'' is itself a claim this paper is obliged to make
precisely. The distance from Theorem~\ref{thm:master} to
$h_{\mathrm{Engine}}$ is \S\S6--7 of \cite{HB1983}; within it the famous
Lemma~10 cites exactly one result from outside the section---Estermann's
bound $S(k;u,v) \ll d(k)\,k^{1/2}(k,u,v)^{1/2}$ for the complete
Kloosterman sum \cite{Estermann}---while \S6, the two-variable
Euler-product apparatus, is the larger share of the formalization.

Two character-sum inputs of that kind are consumed there and only one is
proved in the source. The second, stated at \cite[p.~217]{HB1983} without
proof, is now a theorem of the corpus: for every primitive real character
$\chi \bmod q$ and arbitrary $u, u', v, v'$,
\[
\Bigl|\, \sum_{t \bmod q} \chi(ut + u')\,\chi(vt + v') \,\Bigr|
\;\le\; \gcd\bigl(q,\; uv' - vu'\bigr),
\]
with nothing assumed about the shape of $q$ or of $\chi$ beyond primitivity;
the proof runs through a structure theorem for primitive real characters
that the formalization forced us to make explicit, down to the primitivity
of $\chi_4$---which \textsf{mathlib} names in prose as the primitive
quadratic character on $\mathbb{Z}/4$ without proving it.
\leanname{Salt.HB.sum\_two\_forms\_le\_gcd\_of\_isPrimitive,
Salt.HB.structure\_of\_isPrimitive}
Estermann's bound the corpus supplies through its own Weil track, at present
up to a factor $2^{v_2(k)/2}$: unbounded in the modulus $k$ in general, and
constant on the moduli the engine actually consumes. Closing that reduction
is the named sub-campaign---and pricing the gap in a factor rather than in a
citation is exactly what the formal route buys. The completion apparatus the
argument needs is kernel-checked already: the sawtooth majorant's Fourier
expansion, for every real $\theta$ and every $K \ge 2$, with its
coefficients bounded termwise by $K/(\pi^2 m^2)$ and in $\ell^1$ by
$6(1 + \log K)$.
\leanname{Salt.Weil.sawtoothMajorant\_fourier\_expansion,
Salt.Weil.norm\_majorantCoeff\_le\_sq, Salt.Weil.tsum\_norm\_majorantCoeff\_le}

\begin{table}[h]
\centering\small
\begin{tabular}{lll}
\toprule
Conditional statement & Residual hypothesis & Campaign \\
\midrule
Thm~\ref{thm:dich}, twin horn & $h_{\mathrm{Engine}}$ & \textsc{hb-engine} \\
Thm~\ref{thm:spine} & the MRT door & S8 / MR program \\
Thm~\ref{thm:wall} & the overshoot budget (\texttt{hbudget}) & \S7 \\
\bottomrule
\end{tabular}
\caption{Every conditional in this paper, with its exact named residual.
Nothing else in the paper is conditional.}
\label{tab:conditional}
\end{table}

% ============================================================================
\section*{Acknowledgments}
% ============================================================================
This work was conducted entirely in the author's personal capacity, on
personal time and personal equipment, and used no confidential or
proprietary information of any employer. The views, findings, and
conclusions expressed here are the author's own and do not represent
the views, positions, or endorsement of the author's employer or of any
other organization. The collaboration with Claude (Anthropic's Fable~5
and Opus models) is described in \S\ref{sec:method} and recorded
commit-by-commit in the repository; it likewise implies no position or
endorsement on the part of Anthropic. The author directed and ratified
every design decision and bears sole responsibility for the claims. The
formalization builds on \textsf{mathlib}, and this work stands on the
shoulders of that community.

% ============================================================================
\appendix
\section{The formal-statement index}\label{app:index}
% ============================================================================
% At release: the table below is generated from the All.lean audit blocks; every
% row's axiom column verified by #audit_axioms. The prose <-> Lean renderings
% that follow the table were transcribed from the tree at commit 244ba62
% (2026-08-03) by an extractor: each block is the declaration's source lines with
% the proof body cut, so re-running the extractor is the way to refresh them.
% Widths: every listing line is <= 93 characters, which fits the 6.3in measure at
% \footnotesize in a 0.5em-per-character tt font. Do not hand-edit the listings.

\begin{table}[h]
\centering\small
\begin{tabular}{lll}
\toprule
Statement & Declaration & Axioms \\
\midrule
Thm~\ref{thm:litt} & \leaninline{zeta\_zero\_free\_region\_littlewood} & the three \\
Thm~\ref{thm:pow} & \leaninline{zeta\_zero\_free\_region\_pow} & the three \\
Thm~\ref{thm:rep} & \leaninline{dh\_repulsion\_ordered} & the three \\
Def~\ref{def:fulcrum} & \leaninline{FulcrumQualityMin} & --- \\
Thm~\ref{thm:fulcrumreal} & \leaninline{fulcrum\_zero\_real} & the three \\
Thm~\ref{thm:dich} & \leaninline{fulcrum\_dichotomy} & the three \\
Thm~\ref{thm:master} & \leaninline{hb\_l2c\_master\_unconditional} & the three \\
Thm~\ref{thm:wall} & \leaninline{neutrality\_rate} & the three \\
Thm~\ref{thm:liouville} & \leaninline{S2Gen\_lamR\_eq\_S1} & the three \\
Thm~\ref{thm:gap} & \leaninline{sufficient\_true\_not\_parityInv} & the three \\
Thm~\ref{thm:ziff} & \leaninline{Z\_implies\_TPC / TPC\_implies\_Z} & the three \\
Thm~\ref{thm:spine} & \leaninline{log\_chowla\_two\_door\_only} & the three \\
\bottomrule
\end{tabular}
\caption{``The three'' $=$ \leaninline{propext, Classical.choice,
Quot.sound}. Full index generated at release.}
\end{table}

\subsection{Reading convention}\label{app:conv}

Each entry below pairs a numbered statement of \S\S\ref{sec:banners}--\ref{sec:spine}
with the declaration that carries it, transcribed from the source at the cited
\leaninline{file:line} with the proof body removed and nothing else changed:
implicit binders, instance binders, coercions and numerals stand as Lean writes
them, and where a section-level \leaninline{variable} supplies a binder the entry
says so. The definitions a statement depends on follow it, so an entry closes
over its own vocabulary; anything dropped for width is named and located in the
entry's line, never silently. Every declaration here is listed in an
\leaninline{\#audit\_axioms} block, which is a build-time gate and not a print
(\leaninline{Salt/Tactic/AuditAxioms.lean:81}): it aborts elaboration, naming the
offender, if a listed declaration's axiom closure leaves
$\{\texttt{propext}, \texttt{Classical.choice}, \texttt{Quot.sound}\}$ or if the
name does not resolve --- so \leaninline{lake build} succeeding is the audit
passing, and each entry cites the manifest line that runs it. Names are stable
across the corpus; line numbers are as of commit \leaninline{244ba62}
(2026-08-03), and the source is the arbiter wherever this rendering and the
declaration could be read apart.

\subsection{The four banners (\S\ref{sec:banners})}

\medskip\noindent\textbf{Theorem~\ref{thm:litt}} (Littlewood-strength zero-free
region). Source \leaninline{Salt/Vk/Littlewood.lean:409}; audit
\leaninline{Salt/Vk/All.lean:52}. The width $\log\log|t|/\log|t|$ is the content;
the witnesses $c = 1/88214$ and $T_0 = \exp\exp(\log C + 400)$ are supplied inside
the proof and are not part of the statement.
{\footnotesize\begin{alltt}
theorem zeta_zero_free_region_littlewood :
    ∃ c T₀ : ℝ, 0 < c ∧ 3 ≤ T₀ ∧ ∀ ρ : ℂ, riemannZeta ρ = 0 → T₀ ≤ |ρ.im| →
      ρ.re ≤ 1 - c * (Real.log (Real.log |ρ.im|) / Real.log |ρ.im|)
\end{alltt}}

\medskip\noindent\textbf{Theorem~\ref{thm:pow}} (the power-saving region,
$\theta = 3/4$). Source \leaninline{Salt/Vk/GrowthPow.lean:1044}; audit
\leaninline{Salt/Vk/All.lean:52}. The paper's $(\log\log|t|)^{-O(1)}$ is pinned in
the declaration at the exponent $3$.
{\footnotesize\begin{alltt}
theorem zeta_zero_free_region_pow :
    ∃ c T₀ : ℝ, 0 < c ∧ 3 ≤ T₀ ∧ ∀ ρ : ℂ, riemannZeta ρ = 0 → T₀ ≤ |ρ.im| →
      ρ.re ≤ 1 - c / ((Real.log |ρ.im|) ^ ((3 : ℝ) / 4)
        * (Real.log (Real.log |ρ.im|)) ^ (3 : ℕ))
\end{alltt}}

\medskip\noindent\textbf{Theorem~\ref{thm:rep}} (Deuring--Heilbronn repulsion,
contract form). Source \leaninline{Salt/SW/TBalR8.lean:1752}; audit
\leaninline{Salt/SW/All.lean:124}. The three constants are bound existentially in
the statement; the witnesses are supplied at \leaninline{Salt/SW/TBalR8.lean:1836}
as $\langle 680, c, 14, \ldots \rangle$, with $c$ the explicit minimum set at
\leaninline{:1776}. The paper's $b = 680$ and $k = 14$ are therefore proof data,
not statement data, and a reader checking the display of \S\ref{sec:banners}
against the declaration should read them there.
{\footnotesize\begin{alltt}
theorem dh_repulsion_ordered : ∃ b c k : ℝ, 0 < b ∧ 0 < c ∧ 0 ≤ k ∧
    ∀ (q : ℕ) [NeZero q] (χ : DirichletCharacter ℂ q),
      χ.IsPrimitive → χ ≠ 1 → χ ^ 2 = 1 → 2 ≤ q →
      ∀ β₀ : ℝ, DirichletCharacter.LFunction χ (β₀ : ℂ) = 0 → 1 / 2 < β₀ → β₀ < 1 →
      ∀ ρ : ℂ, DirichletCharacter.LFunction χ ρ = 0 → ρ.im ≠ 0 → |ρ.im| ≤ 1 →
        16 / 17 ≤ ρ.re → ρ.re < 1 → ρ.re ≤ β₀ →
        (1 - β₀) ≥ c * ((q : ℝ) * (|ρ.im| + 2)) ^ (-(b * (1 - ρ.re)))
          / (Real.log ((q : ℝ) * (|ρ.im| + 2)) + 2) ^ k
\end{alltt}}

\medskip\noindent\textbf{Theorem~\ref{thm:vmvt}} (Vinogradov mean value theorem).
Source \leaninline{Salt/Vmvt/Summit2.lean:151}; audit
\leaninline{Salt/Vmvt/All.lean:39}. The conclusion is the named predicate; its
unfolding and the two constants follow, with the exponent $E(k,r)$ the quantity
the source fixes and the constant $D(k,r)$ free.
{\footnotesize\begin{alltt}
theorem vmvt (k r x : ℕ) (hk : 2 ≤ k) (hr : 1 ≤ r) (hx : 1 ≤ x) : VmvtBound k r x

def VmvtBound (k r x : ℕ) : Prop :=
  (JkI k (k * r) x : ℝ) ≤ vmvtConst k r * (x : ℝ) ^ (vmvtExp k r)

noncomputable def JkI (k b x : ℕ) : ℕ := Jk k b (Finset.Ioc (0 : ℤ) (x : ℤ))

noncomputable def vmvtExp (k r : ℕ) : ℝ :=
  2 * (r : ℝ) * (k : ℝ) - (1 / 2) * (k : ℝ) * ((k : ℝ) + 1) + vmvtEta k r

noncomputable def vmvtEta (k r : ℕ) : ℝ := (1 / 2) * (k : ℝ) ^ 2 * (1 - 1 / (k : ℝ)) ^ r

noncomputable def vmvtConst (k r : ℕ) : ℝ := (vmvtC0 k) ^ r

noncomputable def vmvtC0 (k : ℕ) : ℝ := (k : ℝ) ^ (24 * k ^ 2)
\end{alltt}}

\subsection{The fulcrum (\S\ref{sec:fulcrum})}

\medskip\noindent\textbf{Definition~\ref{def:fulcrum}} (the fulcrum). Source
\leaninline{Salt/Fulcrum/Basic.lean:61}. What is formalized is the family
\leaninline{FulcrumQualityMin C}; $C^\star$ is not a Lean constant, and no
declaration in the corpus asserts the minimality claimed for it in
\S\ref{sec:fulcrum} --- that minimality is a property of the demand audit which
produced the definition (the Fulcrum Hunt, recorded in the repository's
exploration ledger), and is to be read as such. The four things the definition
does \emph{not} require of $\rho$ --- reality, $\mathrm{Re}\,\rho < 1$,
$\mathrm{Re}\,\rho > 1/2$, maximality --- are visible as the absence of those
conjuncts.
{\footnotesize\begin{alltt}
def FulcrumQualityMin (C : ℝ) : Prop :=
  ∀ Q : ℕ, ∃ (q : ℕ) (_ : NeZero q) (χ : DirichletCharacter ℂ q) (ρ : ℂ),
    Q < q ∧ χ.IsPrimitive ∧ χ ^ 2 = 1 ∧ χ ≠ 1 ∧
    LFunction χ ρ = 0 ∧ ‖(1 : ℂ) - ρ‖ * (C * Real.log q) ≤ 1
\end{alltt}}

\medskip\noindent\textbf{Theorem~\ref{thm:fulcrumreal}} (reality is derived).
Source \leaninline{Salt/Fulcrum/Basic.lean:93}; audit
\leaninline{Salt/Fulcrum/All.lean:54}. The zero-free region enters as the
hypothesis \leaninline{hZFR} and the quality threshold as \leaninline{hC}; the
conclusion is the conjunction
$\mathrm{Im}\,\rho = 0 \wedge 1/2 \le \mathrm{Re}\,\rho \wedge
\mathrm{Re}\,\rho < 1$.
{\footnotesize\begin{alltt}
theorem fulcrum_zero_real
    (C c₀ : ℝ) (hc₀pos : 0 < c₀) (hc₀le : c₀ ≤ 1) (hC : 2 ≤ C * c₀)
    (hZFR : ∀ (q : ℕ) [NeZero q] (χ : DirichletCharacter ℂ q), χ.IsPrimitive →
      χ ≠ 1 → ∀ \{ρ : ℂ\}, LFunction χ ρ = 0 → 1 / 2 ≤ ρ.re →
        (χ ^ 2 ≠ 1 ∨ ρ.im ≠ 0) →
        ρ.re ≤ 1 - c₀ / Real.log ((q : ℝ) * (|ρ.im| + 2)))
    \{q : ℕ\} [NeZero q] \{χ : DirichletCharacter ℂ q\} \{ρ : ℂ\}
    (hq3 : 3 ≤ q) (hprim : χ.IsPrimitive) (hne : χ ≠ 1)
    (hzero : LFunction χ ρ = 0)
    (hball : ‖(1 : ℂ) - ρ‖ * (C * Real.log q) ≤ 1) :
    ρ.im = 0 ∧ 1 / 2 ≤ ρ.re ∧ ρ.re < 1
\end{alltt}}

\noindent The numeral instance, with the region supplied from the corpus at
$c_0 = 1/126848$ and the threshold read off as $253696 = 2/c_0$ --- the arm of
$C^\star$ that Remark~\ref{rem:arms} reports. Sources
\leaninline{Salt/Fulcrum/CZeroNumeral.lean:427} and \leaninline{:387}; audit
\leaninline{Salt/Fulcrum/All.lean:54}.
{\footnotesize\begin{alltt}
theorem fulcrum_zero_real_numeral (C : ℝ) (hC : 253696 ≤ C)
    \{q : ℕ\} [NeZero q] \{χ : DirichletCharacter ℂ q\} \{ρ : ℂ\}
    (hq3 : 3 ≤ q) (hprim : χ.IsPrimitive) (hne : χ ≠ 1)
    (hzero : LFunction χ ρ = 0)
    (hball : ‖(1 : ℂ) - ρ‖ * (C * Real.log q) ≤ 1) :
    ρ.im = 0 ∧ 1 / 2 ≤ ρ.re ∧ ρ.re < 1

theorem zero_free_region_all_numeral :
    ∀ (q : ℕ) [NeZero q] (χ : DirichletCharacter ℂ q), χ.IsPrimitive →
      χ ≠ 1 → ∀ \{ρ : ℂ\}, LFunction χ ρ = 0 → 1 / 2 ≤ ρ.re → (χ ^ 2 ≠ 1 ∨ ρ.im ≠ 0) →
        ρ.re ≤ 1 - 1 / 126848 / Real.log ((q : ℝ) * (|ρ.im| + 2))
\end{alltt}}

\subsection{The dichotomy (\S\ref{sec:dichotomy})}

\medskip\noindent The three targets, verbatim: sources
\leaninline{Salt/Basic.lean:25}, \leaninline{Salt/TwinBar/SiegelTwin.lean:76}
and \leaninline{:97}.
{\footnotesize\begin{alltt}
def TwinPrimeConjecture : Prop :=
  ∀ n : ℕ, ∃ p : ℕ, n ≤ p ∧ p.Prime ∧ (p + 2).Prime

def NoSiegelZeros : Prop :=
  ∃ c : ℝ, 0 < c ∧ ∀ (q : ℕ) [NeZero q] (χ : DirichletCharacter ℂ q),
    1 < q → χ.IsPrimitive → χ ^ 2 = 1 → χ ≠ 1 →
    ∀ β : ℝ, LFunction χ β = 0 → β < 1 →
      β ≤ 1 - c / Real.log q

def HeathBrownDichotomy : Prop :=
  TwinPrimeConjecture ∨ NoSiegelZeros
\end{alltt}}

\medskip\noindent\textbf{Theorem~\ref{thm:dich}} (the dichotomy). Sources
\leaninline{Salt/Fulcrum/Dichotomy.lean:82} (the unconditional horn) and
\leaninline{:102} (the composite); audit \leaninline{Salt/Fulcrum/All.lean:54}.
The conditionality of the twin horn is a binder: \leaninline{hEngine} is an
argument of the theorem, so no instance of \leaninline{fulcrum\_dichotomy} exists
without one, and Table~\ref{tab:conditional}'s first row is that binder.
{\footnotesize\begin{alltt}
theorem not_fulcrum_implies_noSiegelZeros \{C : ℝ\} (hC : 0 < C)
    (hnF : ¬ FulcrumQualityMin C) : NoSiegelZeros

theorem fulcrum_dichotomy \{C : ℝ\} (hC : 0 < C)
    (hEngine : FulcrumQualityMin C → TwinPrimeConjecture) :
    HeathBrownDichotomy
\end{alltt}}

\subsection{The keystone (\S\ref{sec:keystone})}

\medskip\noindent\textbf{Theorem~\ref{thm:master}} (the master estimate). Source
\leaninline{Salt/HB/L2cMasterUncond.lean:85}; audit
\leaninline{Salt/HB/All.lean:54}. The modulus is supplied by the section
\leaninline{variable} at \leaninline{Salt/HB/L2cMasterUncond.lean:41}. The
hypothesis packet is four items and nothing else --- \leaninline{hsq},
\leaninline{hz100}, \leaninline{hz8}, \leaninline{hzx} --- which is the sense in
which \S\ref{sec:keystone} calls the estimate unconditional; the paper's $2^{31}$
is \leaninline{L2cCmain}.
{\footnotesize\begin{alltt}
theorem hb_l2c_master_unconditional (χ : DirichletCharacter ℂ q) (hsq : χ ^ 2 = 1) \{z x : ℕ\}
    (hz100 : 100 ^ 16 ≤ z) (hz8 : Lwin x ^ 8 ≤ z) (hzx : (z : ℝ) ^ 3 ≤ x) :
    S2 χ (l2cWindow χ z x) - S1 (l2cWindow χ z x)
      ≤ L2cCmain * ((x : ℝ) / z0 z x)
        + L2cCmain * ((x : ℝ) / Real.log x) * Real.exp (5 * z0 z x)
            * PretenseSum χ (2 * x + 2)
        + L2cCmain * Real.exp (2 * z0 z x)
            * ((x : ℝ) / (z : ℝ) ^ (1 / 8 : ℝ) + (x : ℝ) ^ ((9 : ℝ) / 10)) * Lwin x ^ 3
\end{alltt}}

\noindent The vocabulary of that statement, in order of appearance:
\leaninline{Salt/HB/L2cMaster.lean:356}, \leaninline{Salt/HB/L2cCore.lean:159},
\leaninline{:54}, \leaninline{:57}, \leaninline{Salt/HB/Transfer.lean:176} and
\leaninline{:179}. \leaninline{LamTilde} and \leaninline{excPrimorial} are the
character-indexed originals of the two generic definitions rendered in
\S\ref{app:wall}.
{\footnotesize\begin{alltt}
noncomputable def L2cCmain : ℝ := 2 ^ 31

noncomputable def l2cWindow (χ : DirichletCharacter ℂ q) (z x : ℕ) : Finset ℕ :=
  (Finset.Ioc x (2 * x)).filter (fun n => Nat.Coprime (n * (n + 2)) (q * excPrimorial χ z))

noncomputable def Lwin (x : ℕ) : ℝ := Real.log (2 * (x : ℝ) + 2)

noncomputable def z0 (z x : ℕ) : ℝ := Lwin x / Real.log z

noncomputable def S1 (A : Finset ℕ) : ℝ := ∑ n ∈ A, Λ n * Λ (n + 2)

noncomputable def S2 (χ : DirichletCharacter ℂ q) (A : Finset ℕ) : ℝ :=
  ∑ n ∈ A, LamTilde χ n * LamTilde χ (n + 2)
\end{alltt}}

\subsection{The exchange-rate wall (\S\ref{sec:walls})}\label{app:wall}

\medskip\noindent\textbf{Definition~\ref{def:sign}} (sign functions). Source
\leaninline{Salt/HB/SignChain.lean:39}. Four fields, no modulus.
{\footnotesize\begin{alltt}
structure IsSignFunction (f : ℕ → ℝ) : Prop where
  map_one    : f 1 = 1
  map_mul    : ∀ a b : ℕ, f (a * b) = f a * f b
  abs_le_one : ∀ n : ℕ, |f n| ≤ 1
  prime_pm   : ∀ p : ℕ, p.Prime → f p = 1 ∨ f p = -1
\end{alltt}}

\medskip\noindent\textbf{Theorem~\ref{thm:wall}} (neutrality and the rate).
Sources \leaninline{Salt/HB/SignChain.lean:511} (the unconditional leg),
\leaninline{Salt/HB/SignRate.lean:42} (the budget shape) and \leaninline{:53}
(the rate); audit \leaninline{Salt/HB/All.lean:54}. The split the paper states is
the split in the signature: \leaninline{S1\_le\_S2Gen} carries no hypothesis
beyond \leaninline{hf}, and the upper bound consumes \leaninline{hbudget} ---
Table~\ref{tab:conditional}'s third row.
{\footnotesize\begin{alltt}
theorem S1_le_S2Gen (hf : IsSignFunction f) (A : Finset ℕ) : S1 A ≤ S2Gen f A

noncomputable def EngineBound (Cmain : ℝ) (z x : ℕ) (P : ℝ) : ℝ :=
  Cmain * ((x : ℝ) / z0 z x)
    + Cmain * ((x : ℝ) / Real.log x) * Real.exp (5 * z0 z x) * P
    + Cmain * Real.exp (2 * z0 z x)
        * ((x : ℝ) / (z : ℝ) ^ ((1 : ℝ) / 8) + (x : ℝ) ^ ((9 : ℝ) / 10)) * Lwin x ^ 3

theorem neutrality_rate (f : ℕ → ℝ) (hf : IsSignFunction f) (z x : ℕ) (Cmain : ℝ)
    (hbudget : ∑ n ∈ l2cWindowGen f z x, overshootExactGen f n
        ≤ EngineBound Cmain z x (PretenseSumGen f (2 * x + 2))) :
    0 ≤ S2Gen f (l2cWindowGen f z x) - S1 (l2cWindowGen f z x) ∧
      S2Gen f (l2cWindowGen f z x) - S1 (l2cWindowGen f z x)
        ≤ EngineBound Cmain z x (PretenseSumGen f (2 * x + 2))
\end{alltt}}

\medskip\noindent\textbf{Theorem~\ref{thm:liouville}} (the neutrality point).
Source \leaninline{Salt/HB/SignLiouville.lean:93, 97, 103, 109}; audit
\leaninline{Salt/HB/All.lean:54}. Four statements, four lines;
\leaninline{lamR} is the real-valued Liouville function.
{\footnotesize\begin{alltt}
theorem LamTildeGen_lamR_eq_vonMangoldt (n : ℕ) : LamTildeGen lamR n = Λ n

theorem S2Gen_lamR_eq_S1 (A : Finset ℕ) : S2Gen lamR A = S1 A

lemma overshootExactGen_lamR (n : ℕ) : overshootExactGen lamR n = 0

theorem PretenseSumGen_lamR_eq_zero (N : ℕ) : PretenseSumGen lamR N = 0
\end{alltt}}

\noindent The generic vocabulary, all from \leaninline{Salt/HB/SignChain.lean}
(lines 52, 80, 84, 89, 95, 100). \leaninline{S1} is shared verbatim with
\S\ref{sec:keystone}: the transfer is re-typed on one side of the comparison
only.
{\footnotesize\begin{alltt}
noncomputable def LamTildeGen (f : ℕ → ℝ) (n : ℕ) : ℝ :=
  ∑ d ∈ n.divisors, (μ d : ℝ) ^ 2 * f d * Real.log ((n / d : ℕ) : ℝ)

noncomputable def S2Gen (f : ℕ → ℝ) (A : Finset ℕ) : ℝ :=
  ∑ n ∈ A, LamTildeGen f n * LamTildeGen f (n + 2)

noncomputable def overshootExactGen (f : ℕ → ℝ) (n : ℕ) : ℝ :=
  (LamTildeGen f n - Λ n) * LamTildeGen f (n + 2) + Λ n * (LamTildeGen f (n + 2) - Λ (n + 2))

noncomputable def PretenseSumGen (f : ℕ → ℝ) (N : ℕ) : ℝ :=
  ∑ p ∈ (Finset.range (N + 1)).filter (fun p => Nat.Prime p ∧ f p = 1), Real.log (p : ℝ) / p

noncomputable def excPrimorialGen (f : ℕ → ℝ) (z : ℕ) : ℕ :=
  ∏ p ∈ (Finset.range z).filter (fun p => p.Prime ∧ f p ≠ -1), p

noncomputable def l2cWindowGen (f : ℕ → ℝ) (z x : ℕ) : Finset ℕ :=
  (Finset.Ioc x (2 * x)).filter (fun n => Nat.Coprime (n * (n + 2)) (excPrimorialGen f z))
\end{alltt}}

\subsection{The parity gap (\S\ref{sec:parity})}

\medskip\noindent\textbf{Definition~\ref{def:z}} (completions, parity invariance,
$Z$). Source \leaninline{Salt/Parity/Z.lean:44--103}, in file order. Two features
claimed in \S\ref{sec:parity} are checkable here by inspection:
\leaninline{typeISum} is the only functional through which
\leaninline{Completion} reads the weighting, so no bilinear datum is expressible
in the class; and \leaninline{twinRho} is a fixed function, not a quantified one,
so every completion carries the same main term.
{\footnotesize\begin{alltt}
def twinRho (d : ℕ) : ℕ :=
  ((Finset.range d).filter (fun r => d ∣ r * (r + 2))).card

noncomputable def typeISum (a : ℕ → ℝ) (d x : ℕ) : ℝ :=
  ∑ n ∈ Finset.Icc 1 x, if d ∣ n * (n + 2) then a n else 0

noncomputable def typeIError (a : ℕ → ℝ) (θ : ℝ) (x : ℕ) : ℝ :=
  ∑ d ∈ Finset.Icc 1 ⌊(x : ℝ) ^ θ⌋₊,
    |typeISum a d x - (twinRho d : ℝ) / d * x|

def Completion (θ A₀ : ℝ) (a : ℕ → ℝ) : Prop :=
  (∀ n, 0 ≤ a n) ∧
  ∃ C : ℝ, 0 < C ∧ ∀ x : ℕ, 2 ≤ x →
    typeIError a θ x ≤ C * x / Real.log x ^ A₀

def ParityInv (θ A₀ : ℝ) (E : (ℕ → ℝ) → Prop) : Prop :=
  ∀ a : ℕ → ℝ, Completion θ A₀ a → E a

noncomputable def twinMass (a : ℕ → ℝ) (x : ℕ) : ℝ :=
  ∑ n ∈ Finset.Icc 1 x, if n.Prime ∧ (n + 2).Prime then a n else 0

def TwinSufficient (θ A₀ : ℝ) (E : (ℕ → ℝ) → Prop) : Prop :=
  ∀ a : ℕ → ℝ, Completion θ A₀ a → E a →
    ∀ C : ℝ, ∃ x : ℕ, C < twinMass a x

def oneWeight : ℕ → ℝ := fun _ => 1

noncomputable def twinFree : ℕ → ℝ :=
  fun n => if n.Prime ∧ (n + 2).Prime then 0 else 1

def Z (θ A₀ : ℝ) : Prop :=
  ∃ E : (ℕ → ℝ) → Prop, TwinSufficient θ A₀ E ∧ E oneWeight
\end{alltt}}

\medskip\noindent\textbf{Theorem~\ref{thm:gap}} (the gap theorem). Source
\leaninline{Salt/Parity/Z.lean:670}; audit \leaninline{Salt/Parity/All.lean:22}.
The paper's ``certified $A_0$ range'' is the pair of binders \leaninline{h1} and
\leaninline{h2}: $1 \le A_0 \le 2$.
{\footnotesize\begin{alltt}
theorem sufficient_true_not_parityInv \{θ A₀ : ℝ\} (h0 : 0 < θ) (h : θ < 1 / 2)
    (h1 : 1 ≤ A₀) (h2 : A₀ ≤ 2) \{E\}
    (hs : TwinSufficient θ A₀ E) (ht : E oneWeight) : ¬ ParityInv θ A₀ E
\end{alltt}}

\medskip\noindent\textbf{Theorem~\ref{thm:ziff}} ($Z$ is exactly the demand).
Sources \leaninline{Salt/Parity/Z.lean:687} and \leaninline{:216}; audit
\leaninline{Salt/Parity/All.lean:22}. The forward direction is certified on
$0 < \theta < 1/2$, $0 \le A_0$; the converse carries no grade hypothesis.
{\footnotesize\begin{alltt}
theorem Z_implies_TPC \{θ A₀ : ℝ\} (h0 : 0 < θ) (h : θ < 1 / 2) (hA : 0 ≤ A₀) :
    Z θ A₀ → TwinPrimeConjecture

theorem TPC_implies_Z \{θ A₀ : ℝ\} : TwinPrimeConjecture → Z θ A₀
\end{alltt}}

\medskip\noindent The grade guard of \S\ref{sec:parity}, which states the
degeneracy rather than leaving it to be found. Source
\leaninline{Salt/Parity/Z.lean:125}; audit \leaninline{Salt/Parity/All.lean:22}.
{\footnotesize\begin{alltt}
theorem Z_trivial_of_not_completion \{θ A₀ : ℝ\}
    (h : ¬ Completion θ A₀ oneWeight) : Z θ A₀
\end{alltt}}

\medskip\noindent\textbf{Proposition~\ref{prop:census}} (the census). Source
\leaninline{Salt/Parity/Instances.lean}, ten declarations at lines 37, 45, 53,
62, 72, 79, 86, 93, 99, 106: \leaninline{parityInv\_S1\_le\_S2},
\leaninline{parityInv\_twin\_bar}, \leaninline{parityInv\_twin\_gate\_fails},
\leaninline{parityInv\_no\_twin\_weight}, \leaninline{parityInv\_noSiegel\_iff},
\leaninline{parityInv\_chen\_headline},
\leaninline{parityInv\_twin\_almost\_prime}, \leaninline{parityInv\_N6\_2},
\leaninline{parityInv\_N5\_3}, \leaninline{parityInv\_chen\_second}. Six of the
ten are named in the audit block at \leaninline{Salt/Parity/All.lean:22}; the
four that are not are \leaninline{parityInv\_twin\_gate\_fails},
\leaninline{parityInv\_no\_twin\_weight}, \leaninline{parityInv\_noSiegel\_iff}
and \leaninline{parityInv\_N6\_2}. Each of the ten has the shape below --- the
closed corpus theorem wrapped by \leaninline{parityInv\_of\_closed}
(\leaninline{Salt/Parity/Z.lean:117}) --- so the census is one lemma about where
a closed proposition sits, instantiated ten times.
{\footnotesize\begin{alltt}
theorem parityInv_S1_le_S2 (θ A₀ : ℝ) :
    Salt.Parity.ParityInv θ A₀
      (fun _ => ∀ (q : ℕ) (χ : DirichletCharacter ℂ q), χ ^ 2 = 1 →
        ∀ A : Finset ℕ, Salt.HB.S1 A ≤ Salt.HB.S2 χ A)
\end{alltt}}

\subsection{The door-only spine (\S\ref{sec:spine})}

\medskip\noindent\textbf{Theorem~\ref{thm:spine}} (log-Chowla, door-only).
Sources \leaninline{Salt/Entropy/Chowla/SpineFinal.lean:981} (the terminal) and
\leaninline{:750} (the $\forall$-quantified head); audit
\leaninline{Salt/Entropy/All.lean:384}. The single residual is the door binder ---
Table~\ref{tab:conditional}'s second row --- and the head's
\leaninline{extraFloor} quantifier is what places the regime floor above a
prescribed threshold.
{\footnotesize\begin{alltt}
theorem log_chowla_two_door_only :
    ∃ (δ₀ : ℝ) (R : ChowlaRegime), 0 < δ₀ ∧
      ∀ δ : ℝ, 0 < δ → δ ≤ δ₀ → MRTUniformity R δ →
        ¬ logChowla2Fails R.eps R.x R.ω

theorem log_chowla_two_budget_head :
    ∃ (ε : ℚ) (δ₀ : ℝ), 0 < ε ∧ 0 < δ₀ ∧
      ∀ extraFloor : ℕ, ∃ R : ChowlaRegime, R.eps = ε ∧ extraFloor ≤ R.Hlo ∧
        ∀ δ : ℝ, 0 < δ → δ ≤ δ₀ → MRTUniformityXi R δ →
          ¬ logChowla2Fails R.eps R.x R.ω
\end{alltt}}

\noindent The door, its $\Xi_H$-restriction (the form the terminal consumes
before weakening), and the failure predicate the theorems negate. Sources
\leaninline{Salt/Entropy/Chowla/MRTDoor.lean:48}, \leaninline{:109} and
\leaninline{Salt/Entropy/Chowla/ChowlaFailure.lean:59}. The negated predicate is
Tao's two-point statement at shift $1$: the display of \S\ref{sec:spine} is
\leaninline{logChowla2Fails} with the inequality reversed.
{\footnotesize\begin{alltt}
noncomputable def MRTUniformity (R : ChowlaRegime) (δ : ℝ) : Prop :=
  ∀ H : ℕ, R.Hlo ≤ H → H ≤ R.Hhi → ∀ α : ℝ,
    (∫ n, ‖windowExpSum H n α‖ ∂(logMeasure R.x R.ω)) ≤ δ * (H : ℝ)

noncomputable def MRTUniformityXi (R : ChowlaRegime) (δ : ℝ) : Prop :=
  ∀ H : ℕ, ∀ [NeZero H], R.Hlo ≤ H → H ≤ R.Hhi → ∀ ξ ∈ bigXi R.eps H,
    (∫ n, ‖windowExpSum H n (-(ξ.val : ℝ) / (H : ℝ))‖ ∂(logMeasure R.x R.ω)) ≤ δ * (H : ℝ)

def logChowla2Fails (eps : ℚ) (x ω : ℕ) : Prop :=
  (eps : ℝ) * Real.log (ω : ℝ) <
    |∑ n ∈ Finset.Ioc (x / ω) x,
      (ArithmeticFunction.liouville n : ℝ)
        * (ArithmeticFunction.liouville (n + 1) : ℝ) / (n : ℝ)|
\end{alltt}}

\noindent The regime carrying the scales, data fields only. Source
\leaninline{Salt/Entropy/Chowla/Regime.lean:56}; the structure's remaining $19$
fields (lines 73--141, \leaninline{hx} through \leaninline{hxbig}) are proof
obligations discharged at construction, elided here for width.
{\footnotesize\begin{alltt}
structure ChowlaRegime where
  x : ℕ
  ω : ℕ
  a : ℕ
  eps : ℚ
  Hlo : ℕ
  Hhi : ℕ
  C0 : ℕ
  J : ℕ
\end{alltt}}

\subsection{Inherited infrastructure (\S\ref{sec:banners})}

\medskip\noindent The two unconditional almost-prime statements
\S\ref{sec:banners} cites. Sources
\leaninline{Salt/Fulcrum/ChenCorollary.lean:34} (audit
\leaninline{Salt/Fulcrum/All.lean:54}) and
\leaninline{Salt/BrunLower/TwinInstance.lean:771} (audit
\leaninline{Salt/BrunLower/All.lean:54}); \leaninline{cardFactors} is
\textsf{mathlib}'s $\Omega$.
{\footnotesize\begin{alltt}
theorem chen_omega_prod_le_three :
    \{p : ℕ | p.Prime ∧ cardFactors (p * (p + 2)) ≤ 3\}.Infinite

theorem twin_almost_prime :
    \{n : ℕ | ArithmeticFunction.cardFactors (n * (n + 2)) ≤ 20\}.Infinite
\end{alltt}}

% Ordered alphabetically by first-author surname, cited numerically, per the
% Forum of Mathematics Pi instructions for contributors.
\begin{thebibliography}{99}
\bibitem{Chen1973} J.~R. Chen, \emph{On the representation of a larger
even integer as the sum of a prime and the product of at most two primes},
Sci.\ Sinica \textbf{16} (1973), 157--176.
\bibitem{Estermann} T.~Estermann, \emph{On Kloosterman's sum},
Mathematika \textbf{8} (1961), 83--86.
\bibitem{HalesKepler} T.~Hales et al., \emph{A formal proof of the Kepler
conjecture}, Forum of Mathematics, Pi \textbf{5} (2017), e2.
\bibitem{HB1983} D.~R. Heath-Brown, \emph{Prime twins and Siegel zeros},
Proc.\ London Math.\ Soc.\ (3) \textbf{47} (1983), 193--224.
\bibitem{Littlewood1922} J.~E. Littlewood, \emph{Researches in the theory
of the Riemann $\zeta$-function}, Proc.\ London Math.\ Soc.\ (2)
\textbf{20} (1922).
\bibitem{mathlib} The mathlib community, \emph{The Lean mathematical
library}, CPP 2020.
\bibitem{MR2016} K.~Matom\"aki and M.~Radziwi\l\l, \emph{Multiplicative
functions in short intervals}, Ann.\ of Math.\ \textbf{183} (2016).
\bibitem{MRT} K.~Matom\"aki, M.~Radziwi\l\l, and T.~Tao, \emph{An
averaged form of Chowla's conjecture}, Algebra Number Theory \textbf{9}
(2015).
\bibitem{Maynard2015} J.~Maynard, \emph{Small gaps between primes},
Ann.\ of Math.\ \textbf{181} (2015), 383--413.
\bibitem{TaoChowla} T.~Tao, \emph{The logarithmically averaged Chowla and
Elliott conjectures for two-point correlations}, Forum Math.\ Pi
\textbf{4} (2016).
\end{thebibliography}

\end{document}
